%!TEX root = ../main.tex

\chapter{Prompt Prism per Fase 3 e Fase 4}

Il seguente prompt e pensato per generare in Prism la relazione finale completa, usando come base i materiali delle Fasi 1, 2, 3 e 4.

\section{Prompt da copiare in Prism}

\begin{quote}\small
Scrivi una relazione finale di tirocinio universitario in italiano, con tono accademico, chiaro, ordinato e professionale. La relazione deve avere un corpo centrale di circa dieci pagine ed essere adatta a successiva impaginazione in PDF, Word o LaTeX.

Usa come fonti principali i materiali gia preparati: Fase 1, raccolta informazioni del tirocinio; Fase 2, outline dettagliato della relazione; Fase 3, bozza completa della relazione; Fase 4, integrazione approfondita del ruolo della tutor.

Il progetto si chiama \builder{} ed e un'applicazione web per la progettazione, visualizzazione, traduzione ed esportazione di schemi Entity-Relationship e schemi logici. Il progetto e stato svolto come tirocinio universitario. La tutor del tirocinio e stata la Prof.ssa Alessandra Lumini.

La relazione deve descrivere: contesto del tirocinio; obiettivi del lavoro; descrizione generale di \builder; analisi dei requisiti; tecnologie utilizzate; architettura del sistema; funzionalita implementate o consolidate; test, bug fixing e miglioramenti; competenze acquisite; risultati finali, limiti e sviluppi futuri.

Integra in modo esplicito e professionale il ruolo della Prof.ssa Alessandra Lumini come tutor. Spiega che ha supervisionato il tirocinio, fornito consigli metodologici e suggerimenti progettuali, aiutato a definire le priorita e orientato il lavoro verso chiarezza, correttezza concettuale, leggibilita degli schemi, usabilita dell'interfaccia, qualita del codice e valore didattico. Non inventare citazioni dirette e non attribuire alla tutor frasi testuali non documentate.

Usa le informazioni tecniche verificate nel repository: progetto React e TypeScript con Vite; canvas SVG; esportazione PNG, SVG e JPEG; salvataggio/caricamento \texttt{.ersp}; sincronizzazione ERS; traduzione verso schema logico; generazione SQL; reverse engineering SQL; localizzazione in italiano, inglese e albanese; test unitari, test di integrazione e test end-to-end con Playwright; deploy tramite GitHub Actions/GitHub Pages.

Non inventare dati mancanti. Se mancano periodo del tirocinio, ore, CFU, corso di laurea, struttura ospitante, matricola, anno accademico o link del progetto, usa placeholder chiari come \placeholder{Periodo del tirocinio da completare}, \placeholder{Ore/CFU da completare}, \placeholder{Corso di laurea da completare}.

Evita tono pubblicitario, frasi troppo enfatiche, ripetizioni e dettagli tecnici irrilevanti. Mantieni coerenza terminologica: usa sempre \builder{} per il progetto, Prof.ssa Alessandra Lumini per la tutor, schema ER per il modello concettuale e schema logico o modello logico per la trasformazione relazionale.

Organizza il testo in sezioni ben titolate. Collega sempre gli aspetti tecnici al percorso formativo del tirocinio. Concludi con una valutazione realistica dei risultati raggiunti, dei limiti presenti e dei possibili sviluppi futuri.
\end{quote}
